% Part: history
% Chapter: biographies 
% Section: emmy-noether

\documentclass[../../../include/open-logic-section]{subfiles}

\begin{document}

\olfileid[pt]{his}{bio}{noe}

\olsection{Emmy Noether}

\olphoto{noether-emmy}{Emmy Noether}

Emmy Noether (\textsc{ner}-ter) nasceu em Erlangen, na Alemanha, em
23 de março de 1882, em uma família erudita de classe média alta. Aclamada
como a ``mãe da álgebra moderna'', Noether fez contribuições inovadoras tanto
para a matemática quanto para a física, apesar das significativas barreiras à
educação das mulheres. Na Alemanha daquela época, esperava-se que as meninas
fossem educadas nas artes e elas não tinham permissão para frequentar escolas
preparatórias para a faculdade. No entanto, após assistir como ouvinte às
aulas nas Universidades de G\"{o}ttingen e Erlangen (onde seu pai era
professor de matemática), Noether finalmente conseguiu se matricular como
aluna em Erlangen em 1904, quando a política da instituição foi atualizada
para permitir estudantes do sexo feminino. Ela recebeu seu doutorado em
matemática em 1907.

Apesar de suas qualificações, Noether enfrentou muita resistência durante
sua carreira. De 1908 a 1915, ela lecionou em Erlangen sem remuneração.
Durante esse período, chamou a atenção de David Hilbert, um dos mais
eminentes matemáticos do mundo na época, que a convidou para G\"{o}ttingen.
No entanto, as mulheres eram proibidas de obter cátedras, e ela só pôde
lecionar sob o nome de Hilbert, novamente sem remuneração. Foi nesse período
que ela provou o que hoje é conhecido como o teorema de Noether, que ainda é
usado na física teórica atualmente. Noether finalmente recebeu o direito de
lecionar em 1919. A resposta de Hilbert à resistência contínua de seus
colegas de universidade teria sido: ``Senhores, o conselho da faculdade não é
uma casa de banhos.''

No final da década de 1920, ela concentrou-se no trabalho em álgebra abstrata,
e suas contribuições revolucionaram o campo. Em suas provas, ela frequentemente
fazia uso da chamada condição de cadeia ascendente, que afirma que não há
cadeia infinita estritamente crescente de certos conjuntos. Por
exemplo, certas estruturas algébricas hoje conhecidas como anéis noetherianos
têm a propriedade de que não há sequências infinitas de ideais $I_1
\subsetneq I_2 \subsetneq \dots$. A condição pode ser generalizada para
qualquer ordem parcial (na álgebra, ela diz respeito ao caso especial de
ideais ordenados pela relação de subconjunto), e também podemos considerar a
condição dual de cadeia descendente, na qual toda sequência estritamente
\emph{de}crescente em uma ordem parcial eventualmente termina. Se uma ordem
parcial satisfaz a condição de cadeia descendente, é possível usar
indução ao longo dessa ordem de maneira semelhante àquela em que podemos usar
indução ao longo da ordem $<$ sobre~$\Nat$. Tais ordens são chamadas
\emph{bem fundadas} ou \emph{noetherianas}, e o princípio de prova
correspondente, \emph{indução noetheriana}.

Noether era judia e, quando os nazistas chegaram ao poder em 1933, ela foi
demitida de seu cargo. Felizmente, Noether conseguiu emigrar para os
Estados Unidos para ocupar uma posição temporária em Bryn Mawr, na
Pensilvânia. Durante seu tempo lá, ela também lecionou em Princeton,
embora considerasse a universidade pouco acolhedora para as mulheres
\citep[81]{Dick1981}. Em 1935, Noether passou por uma operação para
remover um tumor uterino. Ela morreu de uma infecção decorrente da
cirurgia e foi enterrada em Bryn Mawr.

\begin{reading} 
Para uma biografia de Noether, veja \citet{Dick1981}. O Perimeter
Institute for Theoretical Physics disponibiliza on-line suas palestras sobre a
vida e a influência de Noether \citep{Perimeter2015}. Se você estiver cansado
de ler, \emph{Stuff You Missed in History Class} tem um podcast sobre a
vida e a influência de Noether \citep{Frey2015}. As obras reunidas de
Noether estão disponíveis no alemão original \citep{Noether1983}.
\end{reading}

\end{document}
